\documentclass[11pt]{article}

\usepackage[a4paper,margin=1in]{geometry}
\usepackage{amsmath,amssymb,amsthm}
\usepackage{graphicx}
\usepackage{booktabs}
\usepackage{hyperref}
\usepackage{cite}

\newtheorem{theorem}{Theorem}
\newtheorem{lemma}[theorem]{Lemma}

\title{A Study of Typesetting with \LaTeX}
\author{First Author \and Second Author \\ Some University}
\date{\today}

\begin{document}

\maketitle

\begin{abstract}
  This paper studies the typographic quality of scientific documents produced
  with \LaTeX{}. We describe our method, report experimental results, and
  discuss the implications for authors who prepare camera-ready manuscripts.

  \textbf{Keywords:} typesetting; document preparation; scientific writing
\end{abstract}

\section{Introduction}
\label{sec:intro}

Typesetting quality directly affects the readability of a paper~\cite{knuth1984tex}.
Our contributions are:
\begin{enumerate}
  \item a reproducible workflow for scientific writing;
  \item an evaluation on several document classes;
  \item an open-source template that readers can reuse.
\end{enumerate}

\section{Related Work}
\label{sec:related}

Document preparation systems have evolved from batch-oriented tools to
interactive editors~\cite{lamport1994latex}.

\section{Method}
\label{sec:method}

Let $\mathcal{X}$ be the input space and $f_\theta$ a parameterised model.
The objective is
\begin{equation}
  \mathcal{L}(\theta) = \frac{1}{N}\sum_{i=1}^{N} \ell\bigl(f_\theta(x_i), y_i\bigr) + \lambda \lVert \theta \rVert_2^2 .
  \label{eq:objective}
\end{equation}

\begin{theorem}
  For every $\varepsilon > 0$ there exists a constant $C$ such that the
  estimator in Equation~\eqref{eq:objective} is consistent.
\end{theorem}

\begin{table}[htbp]
  \centering
  \caption{Comparison of rendering accuracy.}
  \label{tab:result}
  \begin{tabular}{lcc}
    \toprule
    Method & Accuracy & Time \\
    \midrule
    Baseline & 88.2\% & 1.0$\times$ \\
    Ours     & \textbf{91.7\%} & 1.1$\times$ \\
    \bottomrule
  \end{tabular}
\end{table}

Table~\ref{tab:result} reports the comparison; Section~\ref{sec:method}
describes the details.

\section{Conclusion}
\label{sec:conclusion}

We presented a complete workflow for preparing scientific manuscripts.

\bibliographystyle{plain}
\bibliography{refs}

\end{document}
